\documentclass[11pt,a4paper]{article}
\usepackage{curso}
\begin{document}
\portada{Lección}{1.1}{Leyes de la termodinámica y energía libre}

\noindent\textbf{Unidad 1 --- Termodinámica e información.}\quad\textbf{Objetivos:} enunciar las leyes; interpretar
$S=\kB\ln\Omega$; usar $F=U-TS$ y $W\ge\Delta F$. \textbf{Prerrequisitos:} lección 0.3; física básica.

\section{¿Por qué importa?}
Un chip convierte trabajo eléctrico en calor. Para computar \emph{con} la física necesitamos el lenguaje de
la energía, el calor y la entropía; la cantidad central del curso será la energía libre \cite{peliti}.

\begin{intuicion}
Dos cajas de partículas en contacto intercambian energía en choques aleatorios simétricos. Sin dirección
impuesta, las energías medias convergen al mismo valor: el equilibrio. La temperatura es la variable que se iguala.
(La animación usa un modelo cinético: en cada choque, la energía total del par se reparte uniformemente al azar \evS.)
\end{intuicion}

\section{Las leyes}
\begin{itemize}
\item \textbf{Ley cero.} El equilibrio térmico es transitivo; define clases de equivalencia etiquetadas por $T$.
\item \textbf{Primera ley.} $\Delta U=Q+W$, con $Q$ el calor \emph{recibido} y $W$ el trabajo hecho \emph{sobre} el sistema.
\item \textbf{Segunda ley.} Para un sistema aislado $\Delta S\ge0$. Para un sistema en contacto con un baño a $T$
(Clausius): $\dd S\ge\delta Q/T$, equivalente a $\Delta S_{\text{sist}}+\Delta S_{\text{baño}}\ge0$ con $\Delta S_{\text{baño}}=-Q/T$.
\item \textbf{Tercera ley.} $S\to S_{\min}$ (cero para un estado fundamental no degenerado) cuando $T\to0$.
\end{itemize}

\section{Entropía y microestados}
\begin{definicion}[ (entropía de Boltzmann)]
$S=\kB\ln\Omega$, con $\Omega$ el número de microestados compatibles con el macroestado.
\end{definicion}
Para $N$ monedas, $\Omega(n\ \text{caras})=\binom Nn$; con $N=100$, $\Omega(50)=\binom{100}{50}\approx1{,}01\times10^{29}$ y
$\Omega(100)=1$. La fracción de caras tiene media $1/2$ y desviación $\frac{1}{2\sqrt N}$: para $N\sim10^{23}$ las
fluctuaciones relativas son $\sim10^{-12}$. La segunda ley es una afirmación estadística: los macroestados con más
microestados son abrumadoramente más probables.

\section{Energía libre}
A temperatura y volumen fijos, el equilibrio minimiza la energía libre de Helmholtz
\begin{equation}F=U-TS .\end{equation}
\begin{proof}[Justificación]
Sistema más baño a $T$: $\Delta S_{\text{univ}}=\Delta S-Q/T=\Delta S-\Delta U/T$ (sin trabajo) $=-\Delta F/T\ge0$, luego $\Delta F\le0$.
\end{proof}
\textbf{Competencia energía--entropía.} Estado $A$: $E=0$, $\Omega=1$. Estado $B$: $E=\varepsilon$, $\Omega=g$. Entonces
$F_A=0$, $F_B=\varepsilon-\kB T\ln g$ y $\frac{p_B}{p_A}=e^{-(F_B-F_A)/\kB T}=g\,e^{-\varepsilon/\kB T}$. La transición
($F_A=F_B$) ocurre en $T^*=\varepsilon/(\kB\ln g)$.
\begin{ejemplo}
$\varepsilon=0{,}1$ eV, $g=1000$: $T^*=0{,}1/(8{,}617\times10^{-5}\cdot6{,}908)\approx168$ K. A 300 K,
$\Delta F=0{,}1-0{,}02585\cdot6{,}908\approx-0{,}0786$ eV y $p_B=1/(1+e^{\Delta F/\kB T})=1/(1+e^{-3{,}04})\approx0{,}954$.
A 100 K, $p_B\approx0{,}009$.
\end{ejemplo}

\section{Trabajo y energía libre}
\begin{resultado}
En un proceso isotérmico (contacto con un baño a $T$), $W\ge\Delta F$, con igualdad solo en el límite cuasiestático.
$W_{\text{dis}}=W-\Delta F\ge0$ es trabajo disipado como calor.
\end{resultado}
\begin{proof}
$\Delta U=Q+W$ y $\Delta S\ge Q/T\Rightarrow W=\Delta U-Q\ge\Delta U-T\Delta S=\Delta F$.
\end{proof}
En la lección 1.5 esta desigualdad da el costo de Landauer; la simulación de borrado mostró $W(\tau)\approx\kB T\ln2+a/\tau$.
La Unidad 3 la refina con la igualdad de Jarzynski $\langle e^{-\beta W}\rangle=e^{-\beta\Delta F}$.

\begin{aplicacion}
Primera ley: la potencia de un chip es calor a extraer. Segunda ley: borrar exige expulsar entropía. En un TSU, el
ruido térmico lleva al sistema hacia estados de baja $F$: la base del muestreo de Boltzmann (Unidades 4--5).
\end{aplicacion}

\begin{error}
\begin{itemize}[nosep]
\item Calor $\neq$ temperatura. La entropía de un \emph{sistema} puede bajar; la del universo no.
\item «Entropía = desorden» es solo una metáfora. Convenciones de signo de $W$: aquí, trabajo sobre el sistema.
\item El mínimo de $F$ caracteriza el equilibrio a $T$ y $V$ fijos; a $T$ y $p$ fijos se usa $G=H-TS$.
\end{itemize}
\end{error}

\section{Resumen}
Equilibrio y $T$; $\Delta U=Q+W$; $\Delta S_{\text{univ}}\ge0$; $S=\kB\ln\Omega$; mínimo de $F=U-TS$; $W\ge\Delta F$.

\begin{thebibliography}{9}\small
\bibitem[2]{peliti} L. Peliti y S. Pigolotti, \emph{Stochastic Thermodynamics: An Introduction}, Princeton University Press, 2021.
\bibitem[10]{wolpert} D. H. Wolpert, The stochastic thermodynamics of computation, \emph{J. Phys. A} 52, 193001 (2019).
\bibitem[12]{jarz} C. Jarzynski, Nonequilibrium equality for free energy differences, \emph{Phys. Rev. Lett.} 78, 2690 (1997).
\end{thebibliography}
\end{document}
